\documentclass[10pt,a4paper]{article}
\usepackage[utf8]{inputenc}
\usepackage[vietnamese]{babel}
\usepackage{amsmath,amssymb,amsfonts}
\usepackage{geometry}
\geometry{a4paper, left=10mm, right=10mm, top=10mm, bottom=10mm}
\usepackage{tikz}
\usetikzlibrary{shapes,arrows.meta,positioning,calc}
\usepackage[most]{tcolorbox}
\usepackage{array,booktabs,multicol}
\usepackage{fontawesome}
\usepackage{helvet}
\renewcommand{\familydefault}{\sfdefault}
\usepackage{parskip}

\definecolor{primary}{RGB}{112, 48, 160}
\definecolor{teal}{RGB}{0, 128, 128}
\definecolor{orange}{RGB}{230, 108, 44}
\definecolor{green}{RGB}{46, 139, 87}
\definecolor{rose}{RGB}{192, 0, 0}

\pagestyle{empty}

\begin{document}

% --- HEADER BANNER ---
\begin{tcolorbox}[
  enhanced,
  colback=violet!8!white,
  colframe=violet!80!black,
  arc=2.5mm, boxrule=1pt,
  left=3mm, right=3mm, top=2mm, bottom=2mm
]
  \begin{tabular*}{\textwidth}{@{\extracolsep{\fill}} lr}
    {\small\bfseries\color{violet!85!black}\faBook\ SỔ TAY KIẾN THỨC TOÁN 2} & 
    {\small\bfseries\color{teal!85!black}\faGraduationCap\ CHƯƠNG TRÌNH GDPT 2018}
  \end{tabular*}
  \vspace{-1mm}
  \begin{center}
    \rule{0.98\textwidth}{0.6pt}
  \end{center}
  \vspace{-2mm}
  \centering
  {\LARGE\bfseries\color{violet!90!black}\faBookmark\ BÀI 72: ÔN TẬP HÌNH HỌC}\\[1mm]
  {\itshape\footnotesize\color{black!75} Tổng kết điểm, đoạn thẳng, đường thẳng, đường gấp khúc, hình phẳng và hình khối}
\end{tcolorbox}

\vspace{-2mm}

\noindent
\begin{minipage}[t]{0.49\textwidth}\vspace{0pt}
\begin{tcolorbox}[
  enhanced,
  colback=violet!4!white,
  colframe=violet!75!black,
  arc=2mm, boxrule=0.8pt,
  left=2mm, right=2mm, top=1.8mm, bottom=1.8mm,
  title={\footnotesize\bfseries\color{white}\faKey\ 1. Điểm, đoạn thẳng, đường gấp khúc}
]
{\footnotesize
\begin{itemize}\setlength{\itemsep}{1pt}\setlength{\leftskip}{-2.5mm}
  \item \textbf{Điểm và đoạn thẳng:} Đoạn thẳng $AB$ bị giới hạn bởi hai đầu mút $A$ và $B$.
  \item \textbf{Đường thẳng:} Kéo dài vô tận về hai phía.
  \item \textbf{Ba điểm thẳng hàng:} Khi ba điểm cùng nằm trên một đường thẳng.
  \item \textbf{Đường gấp khúc:} Gồm nhiều đoạn thẳng nối tiếp nhau (không thẳng hàng):
  $$\text{Độ dài đ.gấp khúc} = \text{Tổng độ dài các đoạn thẳng}$$
\end{itemize}
}
\end{tcolorbox}
\end{minipage}
\hfill
\begin{minipage}[t]{0.49\textwidth}\vspace{0pt}
\begin{tcolorbox}[
  enhanced,
  colback=teal!4!white,
  colframe=teal!75!black,
  arc=2mm, boxrule=0.8pt,
  left=2mm, right=2mm, top=1.8mm, bottom=1.8mm,
  title={\footnotesize\bfseries\color{white}\faSliders\ 2. Hình phẳng và Hình khối}
]
{\footnotesize
\begin{itemize}\setlength{\itemsep}{1pt}\setlength{\leftskip}{-2.5mm}
  \item \textbf{Hình phẳng:}
  \begin{itemize}\setlength{\itemsep}{0.5pt}\setlength{\leftskip}{-2mm}
    \item Hình tam giác: có $3$ cạnh, $3$ đỉnh.
    \item Hình tứ giác: có $4$ cạnh, $4$ đỉnh (hình chữ nhật, hình vuông là dạng tứ giác đặc biệt).
  \end{itemize}
  \item \textbf{Hình khối trong không gian:}
  \begin{itemize}\setlength{\itemsep}{0.5pt}\setlength{\leftskip}{-2mm}
    \item Khối lập phương: $6$ mặt đều là hình vuông.
    \item Khối hộp chữ nhật: $6$ mặt là hình chữ nhật.
    \item Khối trụ: $2$ mặt đáy tròn phẳng, thân tròn cong.
    \item Khối cầu: tròn đều, lăn được về mọi phía.
  \end{itemize}
\end{itemize}
}
\end{tcolorbox}
\end{minipage}

\vspace{-1.5mm}
\begin{tcolorbox}[
  enhanced,
  colback=orange!4!white,
  colframe=orange!80!black,
  arc=2mm, boxrule=0.8pt,
  left=2mm, right=2mm, top=1.5mm, bottom=1.5mm,
  title={\footnotesize\bfseries\color{white}\faSitemap\ 3. Sơ đồ phân loại hệ thống hình học lớp 2}
]
\centering
\begin{tikzpicture}[x=1.45cm, y=0.75cm]
  \tikzset{arr/.style={-{Stealth[scale=0.8]}}, every node/.append style={align=center}}
  \node[rectangle, rounded corners=2pt, draw=violet!80!black, fill=violet!15, font=\scriptsize\bfseries] (hh) at (0,0) {HÌNH HỌC\\LỚP 2};
  \node[rectangle, rounded corners=2pt, draw=teal!80!black, fill=teal!15, font=\tiny\bfseries] (duong) at (2.2,0.8) {ĐƯỜNG - ĐOẠN\\Điểm, đ.thẳng, đ.gấp khúc};
  \node[rectangle, rounded corners=2pt, draw=orange!80!black, fill=orange!15, font=\tiny\bfseries] (phang) at (2.2,0) {HÌNH PHẲNG\\Tam giác, tứ giác, tròn};
  \node[rectangle, rounded corners=2pt, draw=green!70!black, fill=green!15, font=\tiny\bfseries] (khoi) at (2.2,-0.8) {HÌNH KHỐI\\Lập phương, chữ nhật, trụ, cầu};

  \draw[arr, thick] (hh) |- (duong);
  \draw[arr, thick] (hh) -- (phang);
  \draw[arr, thick] (hh) |- (khoi);
\end{tikzpicture}
\end{tcolorbox}

\vspace{-1.5mm}
\noindent
\begin{minipage}[t]{0.49\textwidth}\vspace{0pt}
\begin{tcolorbox}[
  enhanced,
  colback=green!4!white,
  colframe=green!60!black,
  arc=2mm, boxrule=0.8pt,
  left=2mm, right=2mm, top=1.8mm, bottom=1.8mm,
  title={\footnotesize\bfseries\color{white}\faCheckCircle\ 4. Bài tập thực hành mẫu}
]
{\footnotesize
\textbf{1. Tính độ dài đường gấp khúc:}\\
Đường gấp khúc $ABCD$ có các đoạn: $AB = 12\ cm$, $BC = 15\ cm$, $CD = 18\ cm$.\\
$\implies$ Độ dài đường gấp khúc $ABCD$ là:\\
$12 + 15 + 18 = \mathbf{45\ cm}$.\\
\textbf{2. Đếm hình:}\\
Một hình vuông được kẻ thêm $1$ đường chéo sẽ tạo thành \textbf{$2$ hình tam giác} và \textbf{$1$ hình tứ giác} (chính là hình vuông ban đầu).
}
\end{tcolorbox}
\end{minipage}
\hfill
\begin{minipage}[t]{0.49\textwidth}\vspace{0pt}
\begin{tcolorbox}[
  enhanced,
  colback=rose!4!white,
  colframe=rose!70!black,
  arc=2mm, boxrule=0.8pt,
  left=2mm, right=2mm, top=1.8mm, bottom=1.8mm,
  title={\footnotesize\bfseries\color{white}\faExclamationTriangle\ 5. Lỗi đếm hình thiếu}
]
{\footnotesize
\begin{itemize}\setlength{\itemsep}{1.5pt}\setlength{\leftskip}{-2.5mm}
  \item \textbf{Đếm thiếu hình ghép:} Khi đếm tam giác hay tứ giác, học sinh thường chỉ đếm hình đơn mà quên mất các hình đơn ghép lại tạo thành hình lớn.
  \item \textbf{Mẹo đếm hình:} Đánh số $1, 2, 3\dots$ vào từng mảnh nhỏ, sau đó đếm hình đơn rồi đếm hình ghép đôi, ghép ba.
  \item Nhớ kèm tên đơn vị khi tính độ dài đường gấp khúc.
\end{itemize}
}
\end{tcolorbox}
\end{minipage}


\end{document}
